\documentclass[12pt,a4paper]{article}

\usepackage[utf8]{inputenc}
\usepackage[T1]{fontenc}
\usepackage[french]{babel}
\usepackage{geometry}
\geometry{a4paper, margin=2cm}
\usepackage{xcolor}
\usepackage{listings}
\usepackage{hyperref}

% Configuration des couleurs pour le code Java
\definecolor{codegreen}{rgb}{0,0.6,0}
\definecolor{codegray}{rgb}{0.5,0.5,0.5}
\definecolor{codepurple}{rgb}{0.58,0,0.82}
\definecolor{backcolour}{rgb}{0.95,0.95,0.92}

\lstdefinestyle{mystyle}{
    backgroundcolor=\color{backcolour},   
    commentstyle=\color{codegreen},
    keywordstyle=\color{magenta},
    numberstyle=\tiny\color{codegray},
    stringstyle=\color{codepurple},
    basicstyle=\ttfamily\footnotesize,
    breakatwhitespace=false,         
    breaklines=true,                 
    captionpos=b,                    
    keepspaces=true,                 
    numbers=left,                    
    numbersep=5pt,                  
    showspaces=false,                
    showstringspaces=false,
    showtabs=false,                  
    tabsize=2
}
\lstset{style=mystyle}

\title{\textbf{Rapport de Débogage et Refactoring Architecural} \\ \large Suppression des \texttt{Optional} et Résolution des Récursions Infinies (Projet JGrapa)}
\author{Charles Baudoux}
\date{\today}

\begin{document}

\maketitle

\begin{abstract}
Ce rapport détaille les étapes de refactoring effectuées sur le module d'agrégation du projet JGrapa. L'objectif principal était de simplifier l'architecture en remplaçant l'utilisation des \texttt{Optional} par des références directes pouvant être nulles, tout en résolvant les erreurs en cascade (NullPointerException, StackOverflowError) induites par ce changement structurel.
\end{abstract}

\tableofcontents
\vspace{1cm}

\section{Introduction}
L'architecture initiale du système d'agrégation (classes \texttt{Aggregator}, \texttt{Weighter}, \texttt{Parametric}, \texttt{Owa}) utilisait \texttt{Optional<Aggregator>} pour représenter l'absence de sous-agrégateur par défaut (\texttt{defaultSub}). Cette approche a été jugée trop lourde. Le refactoring a consisté à remplacer ces boîtes par des objets directs (acceptant \texttt{null}). Ce changement a nécessité une mise à jour en profondeur des constructeurs, des méthodes de conversion JSON, et a mis en lumière un problème architectural de boucle infinie.

% ---------------------------------------------------------
\section{Nettoyage des Constructeurs et de l'Héritage}
La première étape a été de modifier la signature de la classe parente \texttt{Aggregator} et de toutes ses classes filles pour qu'elles acceptent directement un \texttt{Aggregator} au lieu d'un \texttt{Optional}.

\subsection{Avant la modification}
L'utilisation de \texttt{Optional} forçait l'emballage systématique des paramètres et l'utilisation de \texttt{checkNotNull} sur une valeur qui était censée représenter une absence.

\begin{lstlisting}[language=Java, caption=Constructeur initial avec Optional]
protected Aggregator(Map<Criterion, Aggregator> subs, Optional<Aggregator> defaultSub) {
    this.subs = ImmutableMap.copyOf(subs);
    this.defaultSub = checkNotNull(defaultSub);
}

// Appel dans la classe fille (ex: Owa)
super(subs, defaultSub == Weighter.FULL_EQUAL_WEIGHTER ? Optional.empty() : Optional.of(defaultSub));
\end{lstlisting}

\subsection{Après la modification}
La classe accepte désormais la valeur \texttt{null} de manière native et légitime.

\begin{lstlisting}[language=Java, caption=Nouveau constructeur simplifié]
protected Aggregator(Map<Criterion, Aggregator> subs, Aggregator defaultSub) {
    this.subs = ImmutableMap.copyOf(subs);
    this.defaultSub = defaultSub; // checkNotNull retire car null est accepte
}

// Appel dans la classe fille (ex: Owa)
super(subs, defaultSub == Weighter.FULL_EQUAL_WEIGHTER ? null : defaultSub);
\end{lstlisting}

% ---------------------------------------------------------
\section{Retrait des méthodes associées à Optional}
Le changement de type a rendu obsolètes les méthodes \texttt{.ifPresent()} et \texttt{.orElse()} de l'API Stream de Java, causant des erreurs de compilation.

\subsection{Avant la modification}
\begin{lstlisting}[language=Java, caption=Logique de fallback avec Optional]
public Aggregator defaultSub() {
    return defaultSub.orElse(Weighter.FULL_EQUAL_WEIGHTER);
}

// Utilisation dans une boucle
defaultSub.ifPresent(agg -> subAggregators.put(criterion, agg));
\end{lstlisting}

\subsection{Après la modification}
Nous sommes revenus à des vérifications de nullité standard et à l'utilisation d'opérateurs ternaires.

\begin{lstlisting}[language=Java, caption=Remplacement par des tests de nullite]
public Aggregator defaultSub() {
    return defaultSub != null ? defaultSub : Weighter.FULL_EQUAL_WEIGHTER;
}

// Utilisation dans une boucle
if (defaultSub != null) {
    subAggregators.put(criterion, defaultSub);
}
\end{lstlisting}

% ---------------------------------------------------------
\section{Résolution du StackOverflowError (Boucle Infinie)}
La modification du fallback (\texttt{FULL\_EQUAL\_WEIGHTER}) a provoqué un \texttt{StackOverflowError} lors de l'exécution des tests. Ce bug se produisait dans la méthode \texttt{equals} : deux objets \texttt{Weighter} par défaut tentaient de comparer leurs \texttt{defaultSub()}, qui renvoyaient eux-mêmes, créant une récursivité infinie.

\subsection{Avant la modification}
La comparaison classique forçait l'appel à \texttt{equals} sur le sous-agrégateur sans condition d'arrêt si celui-ci pointait vers lui-même.

\begin{lstlisting}[language=Java, caption=Methode equals provoquant la boucle]
@Override
public boolean equals(Object o2) {
    if (!(o2 instanceof Weighter)) return false;
    final Weighter t2 = (Weighter) o2;
    return weights.equals(t2.weights) 
        && subs().equals(t2.subs()) 
        && defaultSub().equals(t2.defaultSub()); // Point de rupture
}
\end{lstlisting}

\subsection{Après la modification}
L'utilisation de \texttt{Objects.equals()} permet un court-circuit de référence (comparaison d'adresses mémoires) avant de plonger dans l'arbre des objets, stoppant ainsi la récursivité.

\begin{lstlisting}[language=Java, caption=Securisation avec Objects.equals]
@Override
public boolean equals(Object o2) {
    if (!(o2 instanceof Weighter)) return false;
    final Weighter t2 = (Weighter) o2;
    return weights.equals(t2.weights) 
        && subs().equals(t2.subs()) 
        && Objects.equals(defaultSub(), t2.defaultSub()); // Coupe la boucle
}
\end{lstlisting}

% ---------------------------------------------------------
\section{Adaptation du Convertisseur JSON et des Tests}
La sérialisation JSON entrait également dans une boucle infinie à cause du \texttt{defaultSub} qui s'auto-référençait. Enfin, les tests unitaires ont dû être adaptés à l'environnement d'exécution (anglais vs français).

\subsection{Avant la modification (Tests Schema)}
L'assertion dépendait de la langue locale du système d'exploitation, rendant le test instable.

\begin{lstlisting}[language=Java, caption=Test dependant de la locale]
assertTrue(e0.getMessage().contains("enumeration"), e0.getMessage());
\end{lstlisting}

\subsection{Après la modification}
Pour garantir la robustesse du test sur n'importe quel environnement (intégration continue, différents systèmes d'exploitation), la \texttt{Locale} a été fixée explicitement sur l'anglais pour la durée du test.

\begin{lstlisting}[language=Java, caption=Test robuste bilingue et securite JSON]
@Test
void testWrongSchema() throws Exception {
    // Force l'environnement pour rendre le test predictible
    java.util.Locale.setDefault(java.util.Locale.ENGLISH);
    
    // ... code de validation ...
    
    // Test strict sur les cles anglaises de la bibliotheque
    assertTrue(e0.getMessage().contains("enumeration"), e0.getMessage());
    assertTrue(e1.getMessage().contains("array expected"), e1.getMessage());
}
\end{lstlisting}

\section{Conclusion}
Le refactoring s'est soldé par la suppression totale de l'overhead lié aux objets \texttt{Optional} dans l'architecture d'agrégation. L'intégrité de l'arbre a été préservée, les boucles infinies de référencement croisé (StackOverflow) ont été identifiées et neutralisées via \texttt{Objects.equals}, et les tests ont été fiabilisés. L'ensemble des 38 tests passe désormais avec succès.

\end{document}